\section{Explicit finite data for the complete \texorpdfstring{$c=4$}{c=4} slice}\label{sec:c4-data}
This section prints a compressed, complete finite certificate for
Theorem~\ref{thm:c4full}.  A reader can reconstruct every row by decoding an offset
word, applying the corrected lift, sorting the three absolute values, and then decoding
the mode and sign strings.  The accompanying programs only regenerate and check the
printed data; they are not an oracle or a logical premise of the theorem.

\subsection{The four period-48 support rules}\label{sec:c4-support-data}
The zero-sum offset vectors occurring in the periodic rules are encoded as follows.
\begingroup\scriptsize\[
\begin{array}{c|ccccccc}
\text{symbol}&0&1&2&3&4&5&6\\\hline
\delta&(-2,0,2)&(-1,-1,2)&(-1,0,1)&(0,-2,2)&(0,-1,1)&(0,0,0)&(1,-2,1)\\
\end{array}\]\endgroup
\begingroup\scriptsize\[
\begin{array}{c|ccccc}
\text{symbol}&7&8&9&A&B\\\hline
\delta&(1,-1,0)&(1,0,-1)&(2,-2,0)&(2,-1,-1)&(2,0,-2)\\
\end{array}\]\endgroup
Every listed symbol has coordinate sum zero and lies in
$\{-2,-1,0,1,2\}^3$.
For $\rho\in\{1,5,7,11\}$ and $s\ge0$, put $m=m_\rho+12s$.  The phase used at step $s$
is $(\phi_1+16s,\phi_2+16s)$ modulo $48$.  The parameter and bulk words are
listed below; a word is indexed from $0$.
\[\begin{array}{c|c|c}
\rho&m_\rho&(\phi_1,\phi_2)\\\hline
1&25&(8,20)\\
5&29&(16,16)\\
7&31&(16,16)\\
11&23&(0,0)\\
\end{array}\]
\begingroup\scriptsize
\begin{longtable}{c@{\quad}c@{\quad}l}
\hline $\rho$&\text{word}&\text{offset string}\\\hline\endfirsthead
\hline $\rho$&\text{word}&\text{offset string}\\\hline\endhead
1&$W_{1}$&\texttt{5343679953469A695343679953469A695343679953469A69}\\
1&$W_{2}$&\texttt{9B433343654333469B433343654333469B43334365433346}\\
5&$W_{1}$&\texttt{53469A695343679953469A695343679953469A6953436799}\\
5&$W_{2}$&\texttt{33469B433343654333469B433343654333469B4333436543}\\
7&$W_{1}$&\texttt{5343679953469A695343679953469A695343679953469A69}\\
7&$W_{2}$&\texttt{3343654333469B433343654333469B433343654333469B43}\\
11&$W_{1}$&\texttt{5343679953469A695343679953469A695343679953469A69}\\
11&$W_{2}$&\texttt{3343654333469B433343654333469B433343654333469B43}\\
\hline\end{longtable}\endgroup
Put $n=4m$ and $\kappa=n/2$.  In the two bulk zones use
\[
\begin{aligned}
\delta(j)&=W_1[(j-\phi_1-16s)\bmod48]&& (21\le j\le\kappa-21),\\
\delta(j)&=W_2[(j-\phi_2-16s)\bmod48]&& (\kappa+21\le j\le n-21).
\end{aligned}
\]
On $L=[1,20]$, $M=[\kappa-20,\kappa+20]$ and $T=[n-20,n]$ use
the following strings, according to $s\bmod2$; indices increase from left to right.
Thus the three boundary strings have lengths $20$, $41$, and $21$, respectively;
the middle-window index $q$ is relative to $\kappa$.
\begingroup\scriptsize
\begin{longtable}{c@{\,}c@{\,}c@{\quad}c@{\quad}l}
\hline $\rho$&$s\bmod2$&\text{window}&\text{indices}&\text{offset string}\\\hline\endfirsthead
\hline $\rho$&$s\bmod2$&\text{window}&\text{indices}&\text{offset string}\\\hline\endhead
1&0&L&1:20&\texttt{95679979943679953469}\\
1&0&M&-20:20&\texttt{9953469A695343679997779A653343654333469B4}\\
1&0&T&-20:0&\texttt{334695436943336513469}\\
1&1&L&1:20&\texttt{56543334613679953469}\\
1&1&M&-20:20&\texttt{6953436799534698343343679533469B433343654}\\
1&1&T&-20:0&\texttt{334695469943336B43236}\\
5&0&L&1:20&\texttt{56543336944679953469}\\
5&0&M&-20:20&\texttt{695343679953469A443343679533469B433343654}\\
5&0&T&-20:0&\texttt{367995433343431634236}\\
5&1&L&1:20&\texttt{5679997977A679953469}\\
5&1&M&-20:20&\texttt{9953469A695343654333469A695343654333469B4}\\
5&1&T&-20:0&\texttt{6979654336975343699A7}\\
7&0&L&1:20&\texttt{503434333469A6953436}\\
7&0&M&-20:20&\texttt{469A6953436799534698343343654333469B43334}\\
7&0&T&-20:0&\texttt{334695433346334603236}\\
7&1&L&1:20&\texttt{797334695369A6953436}\\
7&1&M&-20:20&\texttt{43679953469A695343469953679B4333436543334}\\
7&1&T&-20:0&\texttt{331365433344334346997}\\
11&0&L&1:20&\texttt{79779979979973303436}\\
11&0&M&-20:20&\texttt{469A695343675333469A699797754333469B43334}\\
11&0&T&-20:0&\texttt{334695436943336546997}\\
11&1&L&1:20&\texttt{797973433469A6953436}\\
11&1&M&-20:20&\texttt{43679953469A695343469995469B4333436543334}\\
11&1&T&-20:0&\texttt{334695436943369546997}\\
\hline\end{longtable}\endgroup
For a finite support check, write uniquely $s=s_0+12Q$, where
$s_0\in\{0,\ldots,11\}$ and $Q\ge0$, and put
$m_0=m_\rho+12s_0$ and $n_0=4m_0$, so that $n=n_0+576Q$.
There are $4\cdot12=48$ states $(\rho,s_0)$.  In a bulk word write
$j=j_0+48h$.
After fixing the sign of each lifted coordinate, every absolute value is affine in $(Q,h)$
with $h$-coefficient $\pm72$ or $\pm144$.  Splitting the $72$-families by $h\bmod2$
gives progressions of difference $144$.  For each residue modulo $144$, substitution of
the printed words and boundary strings gives adjacent quotient intervals beginning at $0$
and ending at $12Q+\lfloor(3n_0-u)/144\rfloor$, for
$u\in\{1,\ldots,144\}$.  Hence the lifted absolute values
are exactly $[1,3n]$ in all $48$ states.
\subsection{The 192 balance patterns}\label{sec:c4-balance-data}
The complete $192=4\cdot12\cdot4$ mode-and-sign rows are collected in
Appendix~\ref{sec:c4-balance-data-appendix}.  The main proof uses only the finite-state
specification and the polynomial-in-$Q$ identity verified by those rows; the full hexadecimal
strings are placed in the appendix to keep the main text readable.
\subsection{The six bounded constructions}\label{sec:c4-small-data}
For $m\in\{5,7,11,13,17,19\}$ the following compact data reconstruct every signed entry.
For each such $m$, they produce four groups of $m$ rows, hence four $m\times3$ arrays.
The radius-two offset alphabet is
\begingroup\scriptsize\[
\begin{array}{c|ccccccc}
\text{symbol}&0&1&2&3&4&5&6\\\hline
\delta&(-2,0,2)&(-2,1,1)&(-2,2,0)&(-1,-1,2)&(-1,0,1)&(-1,1,0)&(-1,2,-1)\\
\end{array}\]\endgroup
\begingroup\scriptsize\[
\begin{array}{c|ccccccc}
\text{symbol}&7&8&9&A&B&C&D\\\hline
\delta&(0,-2,2)&(0,-1,1)&(0,0,0)&(0,1,-1)&(0,2,-2)&(1,-2,1)&(1,-1,0)\\
\end{array}\]\endgroup
\begingroup\scriptsize\[
\begin{array}{c|ccccc}
\text{symbol}&E&F&G&H&I\\\hline
\delta&(1,0,-1)&(1,1,-2)&(2,-2,0)&(2,-1,-1)&(2,0,-2)\\
\end{array}\]\endgroup
Every vector in this radius-two alphabet has coordinate sum zero and entries in
$\{-2,-1,0,1,2\}$.  Each displayed word has length exactly $4m$ and gives
$\delta_1,\ldots,\delta_{4m}$ in the corrected Kotzig lift.
\begingroup\scriptsize
\begin{longtable}{c@{\quad}l}
\hline $m$&\text{offset word}\\\hline\endfirsthead\hline $m$&\text{offset word}\\\hline\endhead
5&\texttt{D37CDDDD9C9CGGGGDDG4}\\
7&\texttt{G9C97CEDGD98738CGGDGHCDGIGD0}\\
11&\texttt{G9707478CGGD9CDDDGE9C788777877CGIGDGHCDGIGD0}\\
13&\texttt{GGDGDDGGB778CGDGD77CD978CGDGGDD777IGGGGGD988887CGGD0}\\
17&\texttt{DGDDDDG2CGG87CGDDDA8CDDD988874DDD7C7CDG978777CGGGG94GDG8788CD9C7C973}\\
19&\texttt{D58CG9CDGCGCDGG9CGDGDDGDA4DDGCDGDDDD987CDG907877787C9777777878CGHDDGGGD7H983}\\
\hline\end{longtable}\endgroup
Let $T_j=\{a_j,b_j,c_j\}$, $a_j<b_j<c_j=a_j+b_j$, be the sorted triple
obtained from the word, and order group $r$ as $T_{4t+r+1}$ for
$t=0,\ldots,m-1$ and $r=0,\ldots,3$.
For mode $c,a,b$, respectively, put
\[
(\mu_t,\nu_t,x_t)=(c_t,b_t-a_t,c_t),\quad(a_t,b_t+c_t,a_t),\quad(b_t,a_t+c_t,b_t).
\]
Decode $E,H$ as binary words, retaining leading zeroes and the first $m$ bits,
to obtain signs $\varepsilon_t,\eta_t$ in the above row order, and set
\begin{equation}\label{eq:c4-small-row}
R_{r,t}=\left((\varepsilon_t\mu_t+\eta_t\nu_t)/2,
(\varepsilon_t\mu_t-\eta_t\nu_t)/2,-\varepsilon_t x_t\right).
\end{equation}
The switch column lists the nondefault modes; ``--'' means all modes are $c$.
In all three modes one has $x_t=\mu_t$, so each row has sum zero and the
third-column identity follows from the first signed chain in Lemma~\ref{lem:exact}.
\begingroup\scriptsize
\setlength{\LTleft}{0pt}\setlength{\LTright}{0pt}
\begin{longtable}{@{\extracolsep{\fill}}c@{\quad}c@{\quad}l@{\quad}l@{\quad}l@{}}
\hline $m$&$r$&\text{switches}&$E$&$H$\\\hline\endfirsthead
\hline $m$&$r$&\text{switches}&$E$&$H$\\\hline\endhead
5&0&$2:b,4:a$&\texttt{C8}&\texttt{48}\\
5&1&$1:a,2:a,4:b$&\texttt{88}&\texttt{48}\\
5&2&$2:b$&\texttt{18}&\texttt{18}\\
5&3&$1:a,2:a,3:a,4:b$&\texttt{18}&\texttt{18}\\
7&0&$3:a,4:a,5:b$&\texttt{16}&\texttt{C6}\\
7&1&$2:a,3:a,4:a,5:a,6:b$&\texttt{1E}&\texttt{72}\\
7&2&$2:a,4:b,5:a,6:b$&\texttt{66}&\texttt{DA}\\
7&3&$4:b,5:b,6:a$&\texttt{8E}&\texttt{32}\\
11&0&--&\texttt{436}&\texttt{0DA}\\
11&1&--&\texttt{82E}&\texttt{116}\\
11&2&$3:b$&\texttt{60E}&\texttt{2BA}\\
11&3&$2:b$&\texttt{21E}&\texttt{89A}\\
13&0&$0:b$&\texttt{8878}&\texttt{3058}\\
13&1&--&\texttt{48B8}&\texttt{0A38}\\
13&2&$0:a$&\texttt{BA18}&\texttt{0D38}\\
13&3&$2:b$&\texttt{2478}&\texttt{4C38}\\
17&0&--&\texttt{C20F8}&\texttt{14078}\\
17&1&$0:a$&\texttt{81578}&\texttt{34078}\\
17&2&$0:a$&\texttt{80E78}&\texttt{12338}\\
17&3&$0:a$&\texttt{021F8}&\texttt{02478}\\
19&0&$1:b$&\texttt{6185E}&\texttt{0505E}\\
19&1&--&\texttt{D013E}&\texttt{02D8E}\\
19&2&--&\texttt{3603E}&\texttt{8C02E}\\
19&3&--&\texttt{5903E}&\texttt{0511E}\\
\hline\end{longtable}\endgroup
For every displayed row, decoding gives $\sum\varepsilon_t\mu_t=0$ and
$\sum\eta_t\nu_t=0$.  Formula~\eqref{eq:c4-small-row} therefore has zero row and
column sums, while the support words give each absolute value in $[1,12m]$ exactly once.

